% ============================================================================
% Quantitative Research -- shared preamble for all five slide decks
%
% Every deck starts with
%     \documentclass[aspectratio=169,11pt]{beamer}
%     % ============================================================================
% Quantitative Research -- shared preamble for all five slide decks
%
% Every deck starts with
%     \documentclass[aspectratio=169,11pt]{beamer}
%     % ============================================================================
% Quantitative Research -- shared preamble for all five slide decks
%
% Every deck starts with
%     \documentclass[aspectratio=169,11pt]{beamer}
%     % ============================================================================
% Quantitative Research -- shared preamble for all five slide decks
%
% Every deck starts with
%     \documentclass[aspectratio=169,11pt]{beamer}
%     \input{../common/preamble.tex}
%     \newcommand{\currentpart}{<k>}
% so that layout, colours, boxes and code listings are identical everywhere.
% ============================================================================

% ======================== Packages ========================
\usepackage[utf8]{inputenc}
\usepackage[T1]{fontenc}
\usepackage[english]{babel}
\usepackage{graphicx}
\usepackage{booktabs}
\usepackage{amsmath,amssymb,amsfonts}
\usepackage{mathtools}
\usepackage{xcolor}
\usepackage{colortbl}
\usepackage{array}
\usepackage{multicol}
\usepackage{verbatim}          % provides the comment environment
\usepackage{tikz}
\usepackage{pgfplots}
\pgfplotsset{compat=1.18}
\usetikzlibrary{arrows.meta,positioning,fit,backgrounds,shapes.geometric,calc,decorations.pathreplacing}
\usepackage{listings}
\usepackage{hyperref}
\usepackage{pgfpages}          % notes on a second screen (only if enabled)

% ======================== Speaker notes (talk-tracks) ========================
% The decks compile WITHOUT notes by default. To get a version with the
% talk-tracks, uncomment exactly ONE of the lines below and recompile:
% \setbeameroption{show notes}                          % notes as separate pages
% \setbeameroption{show notes on second screen=right}   % slides left, notes right

% ======================== Math fonts ========================
% Beamer typesets text in sans serif; make \mathrm and \mathbf match it
% (otherwise SE, TP, RSS print in serif and \mathbf falls back to non-bold).
\renewcommand{\bfdefault}{bx}   % CM sans has bold extended (bx) only
\DeclareMathAlphabet{\mathrm}{T1}{cmss}{m}{n}
\DeclareMathAlphabet{\mathbf}{T1}{cmss}{bx}{n}
\SetMathAlphabet{\mathrm}{bold}{T1}{cmss}{bx}{n}
\SetMathAlphabet{\mathbf}{bold}{T1}{cmss}{bx}{n}
\SetSymbolFont{operators}{bold}{OT1}{cmss}{bx}{n}
% Beamer re-declares \mathrm at \begin{document}; re-apply the sans-serif version
\AtBeginDocument{\SetMathAlphabet{\mathrm}{normal}{T1}{cmss}{m}{n}%
                 \SetMathAlphabet{\mathrm}{bold}{T1}{cmss}{bx}{n}}

% ======================== Colour palette (matched to plots / Madrid) ========
\definecolor{hsbiblau}{HTML}{3333B2}
\definecolor{hsbiblaudunkel}{HTML}{1C1C6E}
\definecolor{hsbiorange}{HTML}{E69F00}
\definecolor{hsbigrau}{HTML}{9A9A9A}
\definecolor{hsbirot}{HTML}{C0392B}
\definecolor{hsbigruen}{HTML}{2E8B57}
\definecolor{hsbiviolett}{HTML}{8C3CA0}   % 4th comparison colour (e.g. candidate lines)

% ======================== Theme ========================
\usetheme{Madrid}
\usecolortheme{default}
\setbeamertemplate{navigation symbols}{}
\setbeamertemplate{blocks}[rounded][shadow=false]
\setbeamerfont{section in head/foot}{size=\tiny}

% Horizontal section navigation bar in the headline
\setbeamertemplate{headline}{%
  \leavevmode%
  \hbox{%
    \begin{beamercolorbox}[wd=\paperwidth,ht=2.5ex,dp=1.125ex]{section in head/foot}%
      \scriptsize\insertsectionnavigationhorizontal{\paperwidth}{}{\hskip0pt plus1filll}%
    \end{beamercolorbox}%
  }%
}

% Footline: lecture part, current section (if any) and slide number
\newcommand{\footsection}{%
  \setbox0=\hbox{\insertsection}%
  \ifdim\wd0>0pt\ \textbar\ \insertsection\fi}
\setbeamertemplate{footline}{%
  \leavevmode%
  \hbox{%
  \begin{beamercolorbox}[wd=\paperwidth,ht=2.5ex,dp=1.125ex]{author in head/foot}%
    \hspace*{2em}Part~\currentpart\footsection\hfill\insertframenumber{}/\inserttotalframenumber\hspace*{2em}%
  \end{beamercolorbox}%
  }%
}

% Automatic section overview slide. Decks with many sections define
%   \newcommand{\overviewleft}{1-4}\newcommand{\overviewright}{5-8}
% before \begin{document} to get a two-column overview.
\newcommand{\overviewtoc}[1]{%
  \tableofcontents[currentsection,sectionstyle=show/shaded,
                   subsectionstyle=show/show/hide,sections={#1}]}
\newcommand{\sectionoverview}{%
  \ifdefined\overviewleft
    \begin{columns}[T]
      \column{0.49\textwidth}\overviewtoc{\overviewleft}
      \column{0.49\textwidth}\overviewtoc{\overviewright}
    \end{columns}%
  \else
    \tableofcontents[currentsection,sectionstyle=show/shaded,
                     subsectionstyle=show/show/hide]%
  \fi}
\AtBeginSection[]{%
  \begin{frame}{Overview}
    \scriptsize
    \sectionoverview
  \end{frame}%
}

% ======================== Python listing style ========================
\lstdefinestyle{Pystyle}{
  language=Python,
  basicstyle=\ttfamily\footnotesize,
  keywordstyle=\color{hsbiblau},
  commentstyle=\color{hsbigrau}\itshape,
  stringstyle=\color{teal},
  showstringspaces=false,
  literate={~}{{$\sim$}}1
           {–}{{--}}1 {—}{{---}}1 {ä}{{\"a}}1 {ö}{{\"o}}1 {ü}{{\"u}}1
           {Ä}{{\"A}}1 {Ö}{{\"O}}1 {Ü}{{\"U}}1 {ß}{{\ss}}1 {é}{{\'e}}1,
  upquote=true,
  columns=fixed, basewidth=0.5em
}
% Plain style for printed Python output (fixed columns keep tables aligned)
\lstdefinestyle{Pyout}{
  language={},
  basicstyle=\ttfamily\scriptsize,
  keywordstyle={}, commentstyle={}, stringstyle={},
  columns=fixed, keepspaces=true, basewidth=0.52em,
  showstringspaces=false,
  upquote=true,
  literate={~}{{$\sim$}}1 {–}{{--}}1 {—}{{---}}1 {‘}{{`}}1 {’}{{'}}1,
}
\lstset{style=Pystyle}

% ======================== Custom commands / styles ========================
\newcommand{\kernaussage}[1]{\begin{alertblock}{Key takeaway}#1\end{alertblock}}
\newcommand{\inworten}[1]{{\small\textbf{\textcolor{hsbiblaudunkel}{In words:}}~#1}}
\newcommand{\guidingquestion}[1]{\begin{exampleblock}{Guiding question}#1\end{exampleblock}}
\newcommand{\appendixnote}[1]{\par\vfill\vspace{-1.6ex}\hfill{\scriptsize\itshape\textcolor{hsbiblau}{$\rightarrow$ Appendix: #1}}\par\vspace{-1ex}}
\newcommand{\sourcenote}[1]{\vfill\hfill{\tiny\itshape\color{gray}Source: #1}\par}

% Companion notebooks (Jupyter / Google Colab). Every deck defines
%     \newcommand{\notebookpath}{Part_k_<Topic>/Part_k_<Topic>.ipynb}
% \trynotebook{<notebook section>} prints a small clickable pointer at the
% bottom right of a frame that opens the deck's notebook in Google Colab
% (the same notebook lies next to the .tex file for local Jupyter use).
\newcommand{\colaburl}{https://colab.research.google.com/github/ChrisW09/Quantitative-Research-Methods/blob/main/}
\newcommand{\trynotebook}[1]{\par\vfill\vspace{-1.6ex}\hfill{\scriptsize\href{\colaburl\notebookpath}{\textcolor{hsbiorange!85!black}{$\blacktriangleright$\ \textbf{Try it yourself:} notebook, #1 \ -- open in Colab}}}\par\vspace{-1ex}}

% Reusable TikZ styles for concept diagrams
\tikzset{
  flowbox/.style={rounded corners, draw=hsbiblaudunkel, very thick,
                  fill=hsbiblau!10, align=center, inner sep=5pt, minimum height=8mm},
  flowboxo/.style={rounded corners, draw=hsbiorange!85!black, very thick,
                   fill=hsbiorange!18, align=center, inner sep=5pt, minimum height=8mm},
  flowboxg/.style={rounded corners, draw=hsbigrau!90!black, very thick,
                   fill=hsbigrau!15, align=center, inner sep=5pt, minimum height=8mm},
  flowarr/.style={-{Latex[length=2.6mm,width=2mm]}, line width=1pt, draw=hsbiblaudunkel},
}

% ======================== Course overview (same slide in every deck) ========
% \courseoverview highlights the part given by \currentpart.
\newcommand{\cpart}[2]{%
  \ifnum\currentpart=#1\relax
    \textcolor{hsbiblau}{\textbf{#2}}%
  \else
    #2%
  \fi}
\newcommand{\courseoverview}{%
\begin{frame}{The lecture at a glance}
  \begin{center}
    \small\renewcommand{\arraystretch}{1.15}
    \begin{tabular}{@{}l l >{\raggedright\arraybackslash}p{9.2cm}@{}}
      \toprule
      & \textbf{Part} & \textbf{Content} \\
      \midrule
      \ifnum\currentpart=1 $\blacktriangleright$\fi & \cpart{1}{1 -- Descriptive Statistics}
        & \cpart{1}{Basic concepts, levels of measurement, frequencies, measures of central tendency and dispersion} \\
      \ifnum\currentpart=2 $\blacktriangleright$\fi & \cpart{2}{2 -- Graphical Representations}
        & \cpart{2}{Empirical CDF, histogram, boxplot} \\
      \ifnum\currentpart=3 $\blacktriangleright$\fi & \cpart{3}{3 -- Statistical Inference}
        & \cpart{3}{Normal distribution and CLT, confidence intervals, sample size, hypothesis tests, $p$-values, one-sided tests, types of error} \\
      \ifnum\currentpart=4 $\blacktriangleright$\fi & \cpart{4}{4 -- Regression}
        & \cpart{4}{Correlation, least squares, goodness of fit, inference for regression, multiple regression (appendix: residual analysis)} \\
      \ifnum\currentpart=5 $\blacktriangleright$\fi & \cpart{5}{5 -- Classification}
        & \cpart{5}{Conditional probability, logistic regression, classification rules, confusion matrix, base rates, ROC/AUC, Bayes classifiers} \\
      \bottomrule
    \end{tabular}
  \end{center}
  \begin{center}
    \footnotesize All analyses are carried out in \textbf{Python} (NumPy,
    pandas, SciPy, statsmodels, matplotlib); every deck comes with a Python
    script that reproduces its numbers and a Jupyter notebook with all
    examples and exercises -- also in Google Colab.
  \end{center}
\end{frame}}

% ======================== Title information (shared) ========================
\title{Quantitative Research}
\author{Prof.\ Dr.\ Christoph Weisser}
\institute{}
\date{}   % must stay explicitly empty -- otherwise Beamer prints today's date

\graphicspath{{images/}}

%     \newcommand{\currentpart}{<k>}
% so that layout, colours, boxes and code listings are identical everywhere.
% ============================================================================

% ======================== Packages ========================
\usepackage[utf8]{inputenc}
\usepackage[T1]{fontenc}
\usepackage[english]{babel}
\usepackage{graphicx}
\usepackage{booktabs}
\usepackage{amsmath,amssymb,amsfonts}
\usepackage{mathtools}
\usepackage{xcolor}
\usepackage{colortbl}
\usepackage{array}
\usepackage{multicol}
\usepackage{verbatim}          % provides the comment environment
\usepackage{tikz}
\usepackage{pgfplots}
\pgfplotsset{compat=1.18}
\usetikzlibrary{arrows.meta,positioning,fit,backgrounds,shapes.geometric,calc,decorations.pathreplacing}
\usepackage{listings}
\usepackage{hyperref}
\usepackage{pgfpages}          % notes on a second screen (only if enabled)

% ======================== Speaker notes (talk-tracks) ========================
% The decks compile WITHOUT notes by default. To get a version with the
% talk-tracks, uncomment exactly ONE of the lines below and recompile:
% \setbeameroption{show notes}                          % notes as separate pages
% \setbeameroption{show notes on second screen=right}   % slides left, notes right

% ======================== Math fonts ========================
% Beamer typesets text in sans serif; make \mathrm and \mathbf match it
% (otherwise SE, TP, RSS print in serif and \mathbf falls back to non-bold).
\renewcommand{\bfdefault}{bx}   % CM sans has bold extended (bx) only
\DeclareMathAlphabet{\mathrm}{T1}{cmss}{m}{n}
\DeclareMathAlphabet{\mathbf}{T1}{cmss}{bx}{n}
\SetMathAlphabet{\mathrm}{bold}{T1}{cmss}{bx}{n}
\SetMathAlphabet{\mathbf}{bold}{T1}{cmss}{bx}{n}
\SetSymbolFont{operators}{bold}{OT1}{cmss}{bx}{n}
% Beamer re-declares \mathrm at \begin{document}; re-apply the sans-serif version
\AtBeginDocument{\SetMathAlphabet{\mathrm}{normal}{T1}{cmss}{m}{n}%
                 \SetMathAlphabet{\mathrm}{bold}{T1}{cmss}{bx}{n}}

% ======================== Colour palette (matched to plots / Madrid) ========
\definecolor{hsbiblau}{HTML}{3333B2}
\definecolor{hsbiblaudunkel}{HTML}{1C1C6E}
\definecolor{hsbiorange}{HTML}{E69F00}
\definecolor{hsbigrau}{HTML}{9A9A9A}
\definecolor{hsbirot}{HTML}{C0392B}
\definecolor{hsbigruen}{HTML}{2E8B57}
\definecolor{hsbiviolett}{HTML}{8C3CA0}   % 4th comparison colour (e.g. candidate lines)

% ======================== Theme ========================
\usetheme{Madrid}
\usecolortheme{default}
\setbeamertemplate{navigation symbols}{}
\setbeamertemplate{blocks}[rounded][shadow=false]
\setbeamerfont{section in head/foot}{size=\tiny}

% Horizontal section navigation bar in the headline
\setbeamertemplate{headline}{%
  \leavevmode%
  \hbox{%
    \begin{beamercolorbox}[wd=\paperwidth,ht=2.5ex,dp=1.125ex]{section in head/foot}%
      \scriptsize\insertsectionnavigationhorizontal{\paperwidth}{}{\hskip0pt plus1filll}%
    \end{beamercolorbox}%
  }%
}

% Footline: lecture part, current section (if any) and slide number
\newcommand{\footsection}{%
  \setbox0=\hbox{\insertsection}%
  \ifdim\wd0>0pt\ \textbar\ \insertsection\fi}
\setbeamertemplate{footline}{%
  \leavevmode%
  \hbox{%
  \begin{beamercolorbox}[wd=\paperwidth,ht=2.5ex,dp=1.125ex]{author in head/foot}%
    \hspace*{2em}Part~\currentpart\footsection\hfill\insertframenumber{}/\inserttotalframenumber\hspace*{2em}%
  \end{beamercolorbox}%
  }%
}

% Automatic section overview slide. Decks with many sections define
%   \newcommand{\overviewleft}{1-4}\newcommand{\overviewright}{5-8}
% before \begin{document} to get a two-column overview.
\newcommand{\overviewtoc}[1]{%
  \tableofcontents[currentsection,sectionstyle=show/shaded,
                   subsectionstyle=show/show/hide,sections={#1}]}
\newcommand{\sectionoverview}{%
  \ifdefined\overviewleft
    \begin{columns}[T]
      \column{0.49\textwidth}\overviewtoc{\overviewleft}
      \column{0.49\textwidth}\overviewtoc{\overviewright}
    \end{columns}%
  \else
    \tableofcontents[currentsection,sectionstyle=show/shaded,
                     subsectionstyle=show/show/hide]%
  \fi}
\AtBeginSection[]{%
  \begin{frame}{Overview}
    \scriptsize
    \sectionoverview
  \end{frame}%
}

% ======================== Python listing style ========================
\lstdefinestyle{Pystyle}{
  language=Python,
  basicstyle=\ttfamily\footnotesize,
  keywordstyle=\color{hsbiblau},
  commentstyle=\color{hsbigrau}\itshape,
  stringstyle=\color{teal},
  showstringspaces=false,
  literate={~}{{$\sim$}}1
           {–}{{--}}1 {—}{{---}}1 {ä}{{\"a}}1 {ö}{{\"o}}1 {ü}{{\"u}}1
           {Ä}{{\"A}}1 {Ö}{{\"O}}1 {Ü}{{\"U}}1 {ß}{{\ss}}1 {é}{{\'e}}1,
  upquote=true,
  columns=fixed, basewidth=0.5em
}
% Plain style for printed Python output (fixed columns keep tables aligned)
\lstdefinestyle{Pyout}{
  language={},
  basicstyle=\ttfamily\scriptsize,
  keywordstyle={}, commentstyle={}, stringstyle={},
  columns=fixed, keepspaces=true, basewidth=0.52em,
  showstringspaces=false,
  upquote=true,
  literate={~}{{$\sim$}}1 {–}{{--}}1 {—}{{---}}1 {‘}{{`}}1 {’}{{'}}1,
}
\lstset{style=Pystyle}

% ======================== Custom commands / styles ========================
\newcommand{\kernaussage}[1]{\begin{alertblock}{Key takeaway}#1\end{alertblock}}
\newcommand{\inworten}[1]{{\small\textbf{\textcolor{hsbiblaudunkel}{In words:}}~#1}}
\newcommand{\guidingquestion}[1]{\begin{exampleblock}{Guiding question}#1\end{exampleblock}}
\newcommand{\appendixnote}[1]{\par\vfill\vspace{-1.6ex}\hfill{\scriptsize\itshape\textcolor{hsbiblau}{$\rightarrow$ Appendix: #1}}\par\vspace{-1ex}}
\newcommand{\sourcenote}[1]{\vfill\hfill{\tiny\itshape\color{gray}Source: #1}\par}

% Companion notebooks (Jupyter / Google Colab). Every deck defines
%     \newcommand{\notebookpath}{Part_k_<Topic>/Part_k_<Topic>.ipynb}
% \trynotebook{<notebook section>} prints a small clickable pointer at the
% bottom right of a frame that opens the deck's notebook in Google Colab
% (the same notebook lies next to the .tex file for local Jupyter use).
\newcommand{\colaburl}{https://colab.research.google.com/github/ChrisW09/Quantitative-Research-Methods/blob/main/}
\newcommand{\trynotebook}[1]{\par\vfill\vspace{-1.6ex}\hfill{\scriptsize\href{\colaburl\notebookpath}{\textcolor{hsbiorange!85!black}{$\blacktriangleright$\ \textbf{Try it yourself:} notebook, #1 \ -- open in Colab}}}\par\vspace{-1ex}}

% Reusable TikZ styles for concept diagrams
\tikzset{
  flowbox/.style={rounded corners, draw=hsbiblaudunkel, very thick,
                  fill=hsbiblau!10, align=center, inner sep=5pt, minimum height=8mm},
  flowboxo/.style={rounded corners, draw=hsbiorange!85!black, very thick,
                   fill=hsbiorange!18, align=center, inner sep=5pt, minimum height=8mm},
  flowboxg/.style={rounded corners, draw=hsbigrau!90!black, very thick,
                   fill=hsbigrau!15, align=center, inner sep=5pt, minimum height=8mm},
  flowarr/.style={-{Latex[length=2.6mm,width=2mm]}, line width=1pt, draw=hsbiblaudunkel},
}

% ======================== Course overview (same slide in every deck) ========
% \courseoverview highlights the part given by \currentpart.
\newcommand{\cpart}[2]{%
  \ifnum\currentpart=#1\relax
    \textcolor{hsbiblau}{\textbf{#2}}%
  \else
    #2%
  \fi}
\newcommand{\courseoverview}{%
\begin{frame}{The lecture at a glance}
  \begin{center}
    \small\renewcommand{\arraystretch}{1.15}
    \begin{tabular}{@{}l l >{\raggedright\arraybackslash}p{9.2cm}@{}}
      \toprule
      & \textbf{Part} & \textbf{Content} \\
      \midrule
      \ifnum\currentpart=1 $\blacktriangleright$\fi & \cpart{1}{1 -- Descriptive Statistics}
        & \cpart{1}{Basic concepts, levels of measurement, frequencies, measures of central tendency and dispersion} \\
      \ifnum\currentpart=2 $\blacktriangleright$\fi & \cpart{2}{2 -- Graphical Representations}
        & \cpart{2}{Empirical CDF, histogram, boxplot} \\
      \ifnum\currentpart=3 $\blacktriangleright$\fi & \cpart{3}{3 -- Statistical Inference}
        & \cpart{3}{Normal distribution and CLT, confidence intervals, sample size, hypothesis tests, $p$-values, one-sided tests, types of error} \\
      \ifnum\currentpart=4 $\blacktriangleright$\fi & \cpart{4}{4 -- Regression}
        & \cpart{4}{Correlation, least squares, goodness of fit, inference for regression, multiple regression (appendix: residual analysis)} \\
      \ifnum\currentpart=5 $\blacktriangleright$\fi & \cpart{5}{5 -- Classification}
        & \cpart{5}{Conditional probability, logistic regression, classification rules, confusion matrix, base rates, ROC/AUC, Bayes classifiers} \\
      \bottomrule
    \end{tabular}
  \end{center}
  \begin{center}
    \footnotesize All analyses are carried out in \textbf{Python} (NumPy,
    pandas, SciPy, statsmodels, matplotlib); every deck comes with a Python
    script that reproduces its numbers and a Jupyter notebook with all
    examples and exercises -- also in Google Colab.
  \end{center}
\end{frame}}

% ======================== Title information (shared) ========================
\title{Quantitative Research}
\author{Prof.\ Dr.\ Christoph Weisser}
\institute{}
\date{}   % must stay explicitly empty -- otherwise Beamer prints today's date

\graphicspath{{images/}}

%     \newcommand{\currentpart}{<k>}
% so that layout, colours, boxes and code listings are identical everywhere.
% ============================================================================

% ======================== Packages ========================
\usepackage[utf8]{inputenc}
\usepackage[T1]{fontenc}
\usepackage[english]{babel}
\usepackage{graphicx}
\usepackage{booktabs}
\usepackage{amsmath,amssymb,amsfonts}
\usepackage{mathtools}
\usepackage{xcolor}
\usepackage{colortbl}
\usepackage{array}
\usepackage{multicol}
\usepackage{verbatim}          % provides the comment environment
\usepackage{tikz}
\usepackage{pgfplots}
\pgfplotsset{compat=1.18}
\usetikzlibrary{arrows.meta,positioning,fit,backgrounds,shapes.geometric,calc,decorations.pathreplacing}
\usepackage{listings}
\usepackage{hyperref}
\usepackage{pgfpages}          % notes on a second screen (only if enabled)

% ======================== Speaker notes (talk-tracks) ========================
% The decks compile WITHOUT notes by default. To get a version with the
% talk-tracks, uncomment exactly ONE of the lines below and recompile:
% \setbeameroption{show notes}                          % notes as separate pages
% \setbeameroption{show notes on second screen=right}   % slides left, notes right

% ======================== Math fonts ========================
% Beamer typesets text in sans serif; make \mathrm and \mathbf match it
% (otherwise SE, TP, RSS print in serif and \mathbf falls back to non-bold).
\renewcommand{\bfdefault}{bx}   % CM sans has bold extended (bx) only
\DeclareMathAlphabet{\mathrm}{T1}{cmss}{m}{n}
\DeclareMathAlphabet{\mathbf}{T1}{cmss}{bx}{n}
\SetMathAlphabet{\mathrm}{bold}{T1}{cmss}{bx}{n}
\SetMathAlphabet{\mathbf}{bold}{T1}{cmss}{bx}{n}
\SetSymbolFont{operators}{bold}{OT1}{cmss}{bx}{n}
% Beamer re-declares \mathrm at \begin{document}; re-apply the sans-serif version
\AtBeginDocument{\SetMathAlphabet{\mathrm}{normal}{T1}{cmss}{m}{n}%
                 \SetMathAlphabet{\mathrm}{bold}{T1}{cmss}{bx}{n}}

% ======================== Colour palette (matched to plots / Madrid) ========
\definecolor{hsbiblau}{HTML}{3333B2}
\definecolor{hsbiblaudunkel}{HTML}{1C1C6E}
\definecolor{hsbiorange}{HTML}{E69F00}
\definecolor{hsbigrau}{HTML}{9A9A9A}
\definecolor{hsbirot}{HTML}{C0392B}
\definecolor{hsbigruen}{HTML}{2E8B57}
\definecolor{hsbiviolett}{HTML}{8C3CA0}   % 4th comparison colour (e.g. candidate lines)

% ======================== Theme ========================
\usetheme{Madrid}
\usecolortheme{default}
\setbeamertemplate{navigation symbols}{}
\setbeamertemplate{blocks}[rounded][shadow=false]
\setbeamerfont{section in head/foot}{size=\tiny}

% Horizontal section navigation bar in the headline
\setbeamertemplate{headline}{%
  \leavevmode%
  \hbox{%
    \begin{beamercolorbox}[wd=\paperwidth,ht=2.5ex,dp=1.125ex]{section in head/foot}%
      \scriptsize\insertsectionnavigationhorizontal{\paperwidth}{}{\hskip0pt plus1filll}%
    \end{beamercolorbox}%
  }%
}

% Footline: lecture part, current section (if any) and slide number
\newcommand{\footsection}{%
  \setbox0=\hbox{\insertsection}%
  \ifdim\wd0>0pt\ \textbar\ \insertsection\fi}
\setbeamertemplate{footline}{%
  \leavevmode%
  \hbox{%
  \begin{beamercolorbox}[wd=\paperwidth,ht=2.5ex,dp=1.125ex]{author in head/foot}%
    \hspace*{2em}Part~\currentpart\footsection\hfill\insertframenumber{}/\inserttotalframenumber\hspace*{2em}%
  \end{beamercolorbox}%
  }%
}

% Automatic section overview slide. Decks with many sections define
%   \newcommand{\overviewleft}{1-4}\newcommand{\overviewright}{5-8}
% before \begin{document} to get a two-column overview.
\newcommand{\overviewtoc}[1]{%
  \tableofcontents[currentsection,sectionstyle=show/shaded,
                   subsectionstyle=show/show/hide,sections={#1}]}
\newcommand{\sectionoverview}{%
  \ifdefined\overviewleft
    \begin{columns}[T]
      \column{0.49\textwidth}\overviewtoc{\overviewleft}
      \column{0.49\textwidth}\overviewtoc{\overviewright}
    \end{columns}%
  \else
    \tableofcontents[currentsection,sectionstyle=show/shaded,
                     subsectionstyle=show/show/hide]%
  \fi}
\AtBeginSection[]{%
  \begin{frame}{Overview}
    \scriptsize
    \sectionoverview
  \end{frame}%
}

% ======================== Python listing style ========================
\lstdefinestyle{Pystyle}{
  language=Python,
  basicstyle=\ttfamily\footnotesize,
  keywordstyle=\color{hsbiblau},
  commentstyle=\color{hsbigrau}\itshape,
  stringstyle=\color{teal},
  showstringspaces=false,
  literate={~}{{$\sim$}}1
           {–}{{--}}1 {—}{{---}}1 {ä}{{\"a}}1 {ö}{{\"o}}1 {ü}{{\"u}}1
           {Ä}{{\"A}}1 {Ö}{{\"O}}1 {Ü}{{\"U}}1 {ß}{{\ss}}1 {é}{{\'e}}1,
  upquote=true,
  columns=fixed, basewidth=0.5em
}
% Plain style for printed Python output (fixed columns keep tables aligned)
\lstdefinestyle{Pyout}{
  language={},
  basicstyle=\ttfamily\scriptsize,
  keywordstyle={}, commentstyle={}, stringstyle={},
  columns=fixed, keepspaces=true, basewidth=0.52em,
  showstringspaces=false,
  upquote=true,
  literate={~}{{$\sim$}}1 {–}{{--}}1 {—}{{---}}1 {‘}{{`}}1 {’}{{'}}1,
}
\lstset{style=Pystyle}

% ======================== Custom commands / styles ========================
\newcommand{\kernaussage}[1]{\begin{alertblock}{Key takeaway}#1\end{alertblock}}
\newcommand{\inworten}[1]{{\small\textbf{\textcolor{hsbiblaudunkel}{In words:}}~#1}}
\newcommand{\guidingquestion}[1]{\begin{exampleblock}{Guiding question}#1\end{exampleblock}}
\newcommand{\appendixnote}[1]{\par\vfill\vspace{-1.6ex}\hfill{\scriptsize\itshape\textcolor{hsbiblau}{$\rightarrow$ Appendix: #1}}\par\vspace{-1ex}}
\newcommand{\sourcenote}[1]{\vfill\hfill{\tiny\itshape\color{gray}Source: #1}\par}

% Companion notebooks (Jupyter / Google Colab). Every deck defines
%     \newcommand{\notebookpath}{Part_k_<Topic>/Part_k_<Topic>.ipynb}
% \trynotebook{<notebook section>} prints a small clickable pointer at the
% bottom right of a frame that opens the deck's notebook in Google Colab
% (the same notebook lies next to the .tex file for local Jupyter use).
\newcommand{\colaburl}{https://colab.research.google.com/github/ChrisW09/Quantitative-Research-Methods/blob/main/}
\newcommand{\trynotebook}[1]{\par\vfill\vspace{-1.6ex}\hfill{\scriptsize\href{\colaburl\notebookpath}{\textcolor{hsbiorange!85!black}{$\blacktriangleright$\ \textbf{Try it yourself:} notebook, #1 \ -- open in Colab}}}\par\vspace{-1ex}}

% Reusable TikZ styles for concept diagrams
\tikzset{
  flowbox/.style={rounded corners, draw=hsbiblaudunkel, very thick,
                  fill=hsbiblau!10, align=center, inner sep=5pt, minimum height=8mm},
  flowboxo/.style={rounded corners, draw=hsbiorange!85!black, very thick,
                   fill=hsbiorange!18, align=center, inner sep=5pt, minimum height=8mm},
  flowboxg/.style={rounded corners, draw=hsbigrau!90!black, very thick,
                   fill=hsbigrau!15, align=center, inner sep=5pt, minimum height=8mm},
  flowarr/.style={-{Latex[length=2.6mm,width=2mm]}, line width=1pt, draw=hsbiblaudunkel},
}

% ======================== Course overview (same slide in every deck) ========
% \courseoverview highlights the part given by \currentpart.
\newcommand{\cpart}[2]{%
  \ifnum\currentpart=#1\relax
    \textcolor{hsbiblau}{\textbf{#2}}%
  \else
    #2%
  \fi}
\newcommand{\courseoverview}{%
\begin{frame}{The lecture at a glance}
  \begin{center}
    \small\renewcommand{\arraystretch}{1.15}
    \begin{tabular}{@{}l l >{\raggedright\arraybackslash}p{9.2cm}@{}}
      \toprule
      & \textbf{Part} & \textbf{Content} \\
      \midrule
      \ifnum\currentpart=1 $\blacktriangleright$\fi & \cpart{1}{1 -- Descriptive Statistics}
        & \cpart{1}{Basic concepts, levels of measurement, frequencies, measures of central tendency and dispersion} \\
      \ifnum\currentpart=2 $\blacktriangleright$\fi & \cpart{2}{2 -- Graphical Representations}
        & \cpart{2}{Empirical CDF, histogram, boxplot} \\
      \ifnum\currentpart=3 $\blacktriangleright$\fi & \cpart{3}{3 -- Statistical Inference}
        & \cpart{3}{Normal distribution and CLT, confidence intervals, sample size, hypothesis tests, $p$-values, one-sided tests, types of error} \\
      \ifnum\currentpart=4 $\blacktriangleright$\fi & \cpart{4}{4 -- Regression}
        & \cpart{4}{Correlation, least squares, goodness of fit, inference for regression, multiple regression (appendix: residual analysis)} \\
      \ifnum\currentpart=5 $\blacktriangleright$\fi & \cpart{5}{5 -- Classification}
        & \cpart{5}{Conditional probability, logistic regression, classification rules, confusion matrix, base rates, ROC/AUC, Bayes classifiers} \\
      \bottomrule
    \end{tabular}
  \end{center}
  \begin{center}
    \footnotesize All analyses are carried out in \textbf{Python} (NumPy,
    pandas, SciPy, statsmodels, matplotlib); every deck comes with a Python
    script that reproduces its numbers and a Jupyter notebook with all
    examples and exercises -- also in Google Colab.
  \end{center}
\end{frame}}

% ======================== Title information (shared) ========================
\title{Quantitative Research}
\author{Prof.\ Dr.\ Christoph Weisser}
\institute{}
\date{}   % must stay explicitly empty -- otherwise Beamer prints today's date

\graphicspath{{images/}}

%     \newcommand{\currentpart}{<k>}
% so that layout, colours, boxes and code listings are identical everywhere.
% ============================================================================

% ======================== Packages ========================
\usepackage[utf8]{inputenc}
\usepackage[T1]{fontenc}
\usepackage[english]{babel}
\usepackage{graphicx}
\usepackage{booktabs}
\usepackage{amsmath,amssymb,amsfonts}
\usepackage{mathtools}
\usepackage{xcolor}
\usepackage{colortbl}
\usepackage{array}
\usepackage{multicol}
\usepackage{verbatim}          % provides the comment environment
\usepackage{tikz}
\usepackage{pgfplots}
\pgfplotsset{compat=1.18}
\usetikzlibrary{arrows.meta,positioning,fit,backgrounds,shapes.geometric,calc,decorations.pathreplacing}
\usepackage{listings}
\usepackage{hyperref}
\usepackage{pgfpages}          % notes on a second screen (only if enabled)

% ======================== Speaker notes (talk-tracks) ========================
% The decks compile WITHOUT notes by default. To get a version with the
% talk-tracks, uncomment exactly ONE of the lines below and recompile:
% \setbeameroption{show notes}                          % notes as separate pages
% \setbeameroption{show notes on second screen=right}   % slides left, notes right

% ======================== Math fonts ========================
% Beamer typesets text in sans serif; make \mathrm and \mathbf match it
% (otherwise SE, TP, RSS print in serif and \mathbf falls back to non-bold).
\renewcommand{\bfdefault}{bx}   % CM sans has bold extended (bx) only
\DeclareMathAlphabet{\mathrm}{T1}{cmss}{m}{n}
\DeclareMathAlphabet{\mathbf}{T1}{cmss}{bx}{n}
\SetMathAlphabet{\mathrm}{bold}{T1}{cmss}{bx}{n}
\SetMathAlphabet{\mathbf}{bold}{T1}{cmss}{bx}{n}
\SetSymbolFont{operators}{bold}{OT1}{cmss}{bx}{n}
% Beamer re-declares \mathrm at \begin{document}; re-apply the sans-serif version
\AtBeginDocument{\SetMathAlphabet{\mathrm}{normal}{T1}{cmss}{m}{n}%
                 \SetMathAlphabet{\mathrm}{bold}{T1}{cmss}{bx}{n}}

% ======================== Colour palette (matched to plots / Madrid) ========
\definecolor{hsbiblau}{HTML}{3333B2}
\definecolor{hsbiblaudunkel}{HTML}{1C1C6E}
\definecolor{hsbiorange}{HTML}{E69F00}
\definecolor{hsbigrau}{HTML}{9A9A9A}
\definecolor{hsbirot}{HTML}{C0392B}
\definecolor{hsbigruen}{HTML}{2E8B57}
\definecolor{hsbiviolett}{HTML}{8C3CA0}   % 4th comparison colour (e.g. candidate lines)

% ======================== Theme ========================
\usetheme{Madrid}
\usecolortheme{default}
\setbeamertemplate{navigation symbols}{}
\setbeamertemplate{blocks}[rounded][shadow=false]
\setbeamerfont{section in head/foot}{size=\tiny}

% Horizontal section navigation bar in the headline
\setbeamertemplate{headline}{%
  \leavevmode%
  \hbox{%
    \begin{beamercolorbox}[wd=\paperwidth,ht=2.5ex,dp=1.125ex]{section in head/foot}%
      \scriptsize\insertsectionnavigationhorizontal{\paperwidth}{}{\hskip0pt plus1filll}%
    \end{beamercolorbox}%
  }%
}

% Footline: lecture part, current section (if any) and slide number
\newcommand{\footsection}{%
  \setbox0=\hbox{\insertsection}%
  \ifdim\wd0>0pt\ \textbar\ \insertsection\fi}
\setbeamertemplate{footline}{%
  \leavevmode%
  \hbox{%
  \begin{beamercolorbox}[wd=\paperwidth,ht=2.5ex,dp=1.125ex]{author in head/foot}%
    \hspace*{2em}Part~\currentpart\footsection\hfill\insertframenumber{}/\inserttotalframenumber\hspace*{2em}%
  \end{beamercolorbox}%
  }%
}

% Automatic section overview slide. Decks with many sections define
%   \newcommand{\overviewleft}{1-4}\newcommand{\overviewright}{5-8}
% before \begin{document} to get a two-column overview.
\newcommand{\overviewtoc}[1]{%
  \tableofcontents[currentsection,sectionstyle=show/shaded,
                   subsectionstyle=show/show/hide,sections={#1}]}
\newcommand{\sectionoverview}{%
  \ifdefined\overviewleft
    \begin{columns}[T]
      \column{0.49\textwidth}\overviewtoc{\overviewleft}
      \column{0.49\textwidth}\overviewtoc{\overviewright}
    \end{columns}%
  \else
    \tableofcontents[currentsection,sectionstyle=show/shaded,
                     subsectionstyle=show/show/hide]%
  \fi}
\AtBeginSection[]{%
  \begin{frame}{Overview}
    \scriptsize
    \sectionoverview
  \end{frame}%
}

% ======================== Python listing style ========================
\lstdefinestyle{Pystyle}{
  language=Python,
  basicstyle=\ttfamily\footnotesize,
  keywordstyle=\color{hsbiblau},
  commentstyle=\color{hsbigrau}\itshape,
  stringstyle=\color{teal},
  showstringspaces=false,
  literate={~}{{$\sim$}}1
           {–}{{--}}1 {—}{{---}}1 {ä}{{\"a}}1 {ö}{{\"o}}1 {ü}{{\"u}}1
           {Ä}{{\"A}}1 {Ö}{{\"O}}1 {Ü}{{\"U}}1 {ß}{{\ss}}1 {é}{{\'e}}1,
  upquote=true,
  columns=fixed, basewidth=0.5em
}
% Plain style for printed Python output (fixed columns keep tables aligned)
\lstdefinestyle{Pyout}{
  language={},
  basicstyle=\ttfamily\scriptsize,
  keywordstyle={}, commentstyle={}, stringstyle={},
  columns=fixed, keepspaces=true, basewidth=0.52em,
  showstringspaces=false,
  upquote=true,
  literate={~}{{$\sim$}}1 {–}{{--}}1 {—}{{---}}1 {‘}{{`}}1 {’}{{'}}1,
}
\lstset{style=Pystyle}

% ======================== Custom commands / styles ========================
\newcommand{\kernaussage}[1]{\begin{alertblock}{Key takeaway}#1\end{alertblock}}
\newcommand{\inworten}[1]{{\small\textbf{\textcolor{hsbiblaudunkel}{In words:}}~#1}}
\newcommand{\guidingquestion}[1]{\begin{exampleblock}{Guiding question}#1\end{exampleblock}}
\newcommand{\appendixnote}[1]{\par\vfill\vspace{-1.6ex}\hfill{\scriptsize\itshape\textcolor{hsbiblau}{$\rightarrow$ Appendix: #1}}\par\vspace{-1ex}}
\newcommand{\sourcenote}[1]{\vfill\hfill{\tiny\itshape\color{gray}Source: #1}\par}

% Companion notebooks (Jupyter / Google Colab). Every deck defines
%     \newcommand{\notebookpath}{Part_k_<Topic>/Part_k_<Topic>.ipynb}
% \trynotebook{<notebook section>} prints a small clickable pointer at the
% bottom right of a frame that opens the deck's notebook in Google Colab
% (the same notebook lies next to the .tex file for local Jupyter use).
\newcommand{\colaburl}{https://colab.research.google.com/github/ChrisW09/Quantitative-Research-Methods/blob/main/}
\newcommand{\trynotebook}[1]{\par\vfill\vspace{-1.6ex}\hfill{\scriptsize\href{\colaburl\notebookpath}{\textcolor{hsbiorange!85!black}{$\blacktriangleright$\ \textbf{Try it yourself:} notebook, #1 \ -- open in Colab}}}\par\vspace{-1ex}}

% Reusable TikZ styles for concept diagrams
\tikzset{
  flowbox/.style={rounded corners, draw=hsbiblaudunkel, very thick,
                  fill=hsbiblau!10, align=center, inner sep=5pt, minimum height=8mm},
  flowboxo/.style={rounded corners, draw=hsbiorange!85!black, very thick,
                   fill=hsbiorange!18, align=center, inner sep=5pt, minimum height=8mm},
  flowboxg/.style={rounded corners, draw=hsbigrau!90!black, very thick,
                   fill=hsbigrau!15, align=center, inner sep=5pt, minimum height=8mm},
  flowarr/.style={-{Latex[length=2.6mm,width=2mm]}, line width=1pt, draw=hsbiblaudunkel},
}

% ======================== Course overview (same slide in every deck) ========
% \courseoverview highlights the part given by \currentpart.
\newcommand{\cpart}[2]{%
  \ifnum\currentpart=#1\relax
    \textcolor{hsbiblau}{\textbf{#2}}%
  \else
    #2%
  \fi}
\newcommand{\courseoverview}{%
\begin{frame}{The lecture at a glance}
  \begin{center}
    \small\renewcommand{\arraystretch}{1.15}
    \begin{tabular}{@{}l l >{\raggedright\arraybackslash}p{9.2cm}@{}}
      \toprule
      & \textbf{Part} & \textbf{Content} \\
      \midrule
      \ifnum\currentpart=1 $\blacktriangleright$\fi & \cpart{1}{1 -- Descriptive Statistics}
        & \cpart{1}{Basic concepts, levels of measurement, frequencies, measures of central tendency and dispersion} \\
      \ifnum\currentpart=2 $\blacktriangleright$\fi & \cpart{2}{2 -- Graphical Representations}
        & \cpart{2}{Empirical CDF, histogram, boxplot} \\
      \ifnum\currentpart=3 $\blacktriangleright$\fi & \cpart{3}{3 -- Statistical Inference}
        & \cpart{3}{Normal distribution and CLT, confidence intervals, sample size, hypothesis tests, $p$-values, one-sided tests, types of error} \\
      \ifnum\currentpart=4 $\blacktriangleright$\fi & \cpart{4}{4 -- Regression}
        & \cpart{4}{Correlation, least squares, goodness of fit, inference for regression, multiple regression (appendix: residual analysis)} \\
      \ifnum\currentpart=5 $\blacktriangleright$\fi & \cpart{5}{5 -- Classification}
        & \cpart{5}{Conditional probability, logistic regression, classification rules, confusion matrix, base rates, ROC/AUC, Bayes classifiers} \\
      \bottomrule
    \end{tabular}
  \end{center}
  \begin{center}
    \footnotesize All analyses are carried out in \textbf{Python} (NumPy,
    pandas, SciPy, statsmodels, matplotlib); every deck comes with a Python
    script that reproduces its numbers and a Jupyter notebook with all
    examples and exercises -- also in Google Colab.
  \end{center}
\end{frame}}

% ======================== Title information (shared) ========================
\title{Quantitative Research}
\author{Prof.\ Dr.\ Christoph Weisser}
\institute{}
\date{}   % must stay explicitly empty -- otherwise Beamer prints today's date

\graphicspath{{images/}}
